\documentclass[12pt,a4paper]{article}

% ── Packages ──────────────────────────────────────────────────────────────────
\usepackage[margin=2.5cm]{geometry}
\usepackage{graphicx}
\usepackage{xcolor}
\usepackage{titlesec}
\usepackage{booktabs}
\usepackage{longtable}
\usepackage{array}
\usepackage{enumitem}
\usepackage{fancyhdr}
\usepackage{hyperref}
\usepackage{listings}
\usepackage{multicol}
\usepackage{tcolorbox}
\usepackage{tabularx}
\usepackage{float}
\usepackage{parskip}
\usepackage{microtype}
\usepackage{lmodern}
\usepackage[T1]{fontenc}
\usepackage{colortbl}
\usepackage{makecell}

% ── Colours ───────────────────────────────────────────────────────────────────
\definecolor{primary}{HTML}{1A237E}
\definecolor{accent}{HTML}{0D47A1}
\definecolor{lightblue}{HTML}{E3F2FD}
\definecolor{critical}{HTML}{B71C1C}
\definecolor{high}{HTML}{E65100}
\definecolor{medium}{HTML}{F57F17}
\definecolor{infogreen}{HTML}{1B5E20}
\definecolor{codebg}{HTML}{F5F5F5}
\definecolor{tabhead}{HTML}{1A237E}
\definecolor{rowalt}{HTML}{EEF2FF}

% ── Section styling ───────────────────────────────────────────────────────────
\titleformat{\section}
  {\Large\bfseries\color{primary}}{\thesection}{1em}{}[\titlerule]
\titleformat{\subsection}
  {\large\bfseries\color{accent}}{\thesubsection}{1em}{}
\titleformat{\subsubsection}
  {\normalsize\bfseries\color{accent}}{\thesubsubsection}{1em}{}

% ── Header / Footer ───────────────────────────────────────────────────────────
\pagestyle{fancy}
\fancyhf{}
\rhead{\textcolor{primary}{\textbf{SIEM Project Report}}}
\lhead{\textcolor{accent}{Security Information \& Event Management}}
\rfoot{\textcolor{primary}{Page \thepage}}
\lfoot{\textcolor{accent}{\today}}
\renewcommand{\headrulewidth}{0.4pt}
\renewcommand{\footrulewidth}{0.4pt}

% ── Code listing ──────────────────────────────────────────────────────────────
\lstset{
  backgroundcolor=\color{codebg},
  basicstyle=\ttfamily\footnotesize,
  breaklines=true,
  frame=single,
  rulecolor=\color{accent},
  keywordstyle=\color{primary}\bfseries,
  commentstyle=\color{gray}\itshape,
  stringstyle=\color{infogreen},
  numbers=left,
  numberstyle=\tiny\color{gray},
  xleftmargin=1.5em,
  framexleftmargin=1.5em
}

% ── Coloured boxes ────────────────────────────────────────────────────────────
\tcbuselibrary{skins,breakable}
\newtcolorbox{infobox}[1][]{
  colback=lightblue,colframe=accent,
  fonttitle=\bfseries,title=#1,breakable}
\newtcolorbox{critbox}[1][]{
  colback=red!8,colframe=critical,
  fonttitle=\bfseries\color{critical},title=#1,breakable}
\newtcolorbox{rulebox}[1][]{
  colback=yellow!10,colframe=medium,
  fonttitle=\bfseries\color{medium},title=#1,breakable}

% ── Hyperref ──────────────────────────────────────────────────────────────────
\hypersetup{colorlinks=true,linkcolor=primary,urlcolor=accent,citecolor=accent}

% ── Helpers ───────────────────────────────────────────────────────────────────
\newcommand{\file}[1]{\texttt{\color{accent}#1}}
\newcommand{\field}[1]{\texttt{#1}}

% ══════════════════════════════════════════════════════════════════════════════
\begin{document}

% ── Title Page ────────────────────────────────────────────────────────────────
\begin{titlepage}
  \pagecolor{primary}\color{white}
  \vspace*{2.5cm}
  \begin{center}
    {\fontsize{13}{15}\selectfont\texttt{Software Engineering Project — Final Report}}\\[0.6cm]
    \rule{\linewidth}{0.5pt}\\[1.2cm]
    {\fontsize{38}{44}\selectfont\bfseries SIEM}\\[0.5cm]
    {\fontsize{20}{24}\selectfont Security Information \&\\[0.3cm]Event Management System}\\[1.2cm]
    \rule{\linewidth}{0.5pt}\\[1.5cm]
    {\large
      A real-time threat detection and log analysis platform\\[0.3cm]
      powered by the Elastic Stack (ELK) with automated\\[0.3cm]
      alert generation, anomaly detection rules, and\\[0.3cm]
      an interactive Kibana security dashboard.
    }\\[2.5cm]
    \begin{tabular}{ll}
      \textbf{Team}  & SIEM Project Team \\[0.4cm]
      \textbf{Stack} & Python \textbullet\ Elasticsearch \textbullet\ Logstash
                       \textbullet\ Kibana \textbullet\ Filebeat \textbullet\ Docker \\[0.4cm]
      \textbf{Date}  & \today \\
    \end{tabular}
  \end{center}
  \vfill
  \begin{center}
    \small\textit{Submitted in partial fulfilment of the requirements for the\\
    Software Engineering module}
  \end{center}
\end{titlepage}
\nopagecolor

% ── Table of Contents ─────────────────────────────────────────────────────────
\newpage
\tableofcontents
\newpage

% ══════════════════════════════════════════════════════════════════════════════
\section{Abstract}

Modern organisations generate hundreds of thousands of log events every day.
Without automated correlation and analysis, security incidents go undetected
until the damage is catastrophic. This project delivers a fully operational
\textbf{Security Information and Event Management (SIEM)} system that ingests
SSH authentication logs in real time, parses them through a structured
enrichment pipeline, detects threat patterns using five detection rules, and
surfaces everything through live Kibana dashboards.

\subsection*{Problem Statement}

System administrators have traditionally relied on manual log inspection to
identify unauthorised access attempts. This approach fails at scale: a
brute-force campaign can exhaust credentials in under sixty seconds; a
coordinated botnet attack generates thousands of events per minute; and
lateral movement by an attacker who has already gained initial access leaves
only subtle traces across multiple log files.

A purpose-built SIEM reduces \textbf{Mean Time To Detect (MTTD)} from hours —
or never — to under fifteen seconds by automating correlation and alerting.

\subsection*{Target Users}

\begin{itemize}[leftmargin=2em]
  \item \textbf{Security Admin} — primary user; monitors live dashboards,
        reviews audit trails, and manages the IP blocklist.
  \item \textbf{System Operator} — views event feeds to distinguish normal
        from abnormal activity without modifying any configuration.
  \item \textbf{Analyst / Auditor} — queries historical event data for
        compliance reporting and incident post-mortems.
\end{itemize}

\subsection*{Key Implemented Features}

\begin{enumerate}[leftmargin=2em]
  \item \textbf{Realistic log simulation} — \file{logg\_gen.py} produces
        convincing SSH traffic: fixed recurring actor IPs, random attacker
        pools, burst-attack state machine, and forced alert scenarios every
        10 seconds to guarantee detectable events.
  \item \textbf{Full ELK ingestion pipeline} — Filebeat tails log files and
        ships lines to Logstash, which applies 14 Grok patterns and routes
        enriched documents into daily Elasticsearch indices.
  \item \textbf{Python alert engine} — \file{alerts.py} polls Elasticsearch
        every 10 seconds, evaluates five rule-based detectors, and writes
        structured JSON alerts back into the same pipeline.
  \item \textbf{Automatic IP blocklist} — HIGH and CRITICAL alerts
        auto-populate \file{blocklist.json} with timestamp, reason, and
        severity.
  \item \textbf{Kibana security dashboards} — Two fully configured dashboards
        (SSH Logs Overview and Alerts \& Threats) with 14 visualisation panels
        covering metrics, timelines, charts, and sortable audit tables.
\end{enumerate}

\subsection*{Top Challenges}

\begin{itemize}[leftmargin=2em]
  \item Writing Grok patterns that correctly parse all 14 SSH log variants,
        including services with hyphens (\texttt{systemd-logind}), lines
        with and without PIDs, and PAM multi-token formats.
  \item Closing the feedback loop: alert JSON written to a log file is
        re-ingested by Filebeat through the \emph{same} Logstash pipeline
        into a \emph{separate} index — keeping dashboards unified without
        polluting raw log data.
  \item Reconstructing accurate \field{@timestamp} values from SSH log lines
        that carry no year field, across midnight and year boundaries, in the
        correct timezone.
\end{itemize}

\newpage

% ══════════════════════════════════════════════════════════════════════════════
\section{Requirements Specification}

\subsection{Functional Requirements}

\begin{longtable}{>{\bfseries\small}p{0.8cm} >{\small}p{3.8cm} >{\small}p{8.5cm}}
  \toprule
  \rowcolor{tabhead}
  \textcolor{white}{\textbf{ID}} &
  \textcolor{white}{\textbf{Feature}} &
  \textcolor{white}{\textbf{Description}} \\
  \midrule\endhead
  F1 & Log Collection &
    System ingests SSH auth log lines from a watched file within 5 seconds
    of being written, using Filebeat with a 5-second scan frequency. \\
  \rowcolor{rowalt}
  F2 & Log Parsing &
    Every ingested line is decomposed into structured fields: \field{src\_ip},
    \field{username}, \field{event\_type}, \field{severity}, \field{@timestamp},
    \field{service}, \field{pid}, \field{server}. \\
  F3 & Event Classification &
    Fourteen distinct SSH event types are recognised and labelled. Unknown
    lines are tagged \field{unparsed} without dropping them. \\
  \rowcolor{rowalt}
  F4 & Severity Enrichment &
    Each event is assigned a severity level (info / warning / high / critical)
    based on event type and username. Root logins are always critical. \\
  F5 & Alert Generation &
    Five detection rules fire automatically when configurable thresholds are
    exceeded. Alerts are written as structured JSON for re-ingestion. \\
  \rowcolor{rowalt}
  F6 & Duplicate Suppression &
    A 60-second per-IP cooldown prevents the same alert rule firing repeatedly
    for the same ongoing incident. \\
  F7 & IP Blocklist &
    HIGH and CRITICAL alerts auto-add the attacker IP to a persistent JSON
    blocklist with blocked timestamp and reason. \\
  \rowcolor{rowalt}
  F8 & Admin Dashboard &
    Kibana provides real-time charts, metric cards, time-series graphs, and
    sortable event tables — no coding required to use. \\
  F9 & Audit Trail &
    All events are queryable by time range, IP address, username, and severity,
    providing a full historical record of system activity. \\
  \rowcolor{rowalt}
  F10 & Alert Index &
    Fired alerts are stored in a separate \field{siem-alerts-*} index,
    queryable independently of raw log data. \\
  \bottomrule
  \caption{Functional Requirements}
\end{longtable}

\subsection{Non-Functional Requirements}

\begin{longtable}{>{\bfseries\small}p{0.8cm} >{\small}p{3cm} >{\small}p{9.5cm}}
  \toprule
  \rowcolor{tabhead}
  \textcolor{white}{\textbf{ID}} &
  \textcolor{white}{\textbf{Category}} &
  \textcolor{white}{\textbf{Requirement}} \\
  \midrule\endhead
  NF1 & Performance &
    End-to-end latency from log write to Kibana visibility is under 15 seconds
    under normal load (Filebeat scan: 5s + Logstash processing: $<$1s +
    Elasticsearch indexing: $<$1s). \\
  \rowcolor{rowalt}
  NF2 & Scalability &
    Elasticsearch single-node handles $>$10,000 events/day in development.
    Production scaling via replica shards and Logstash worker threads. \\
  NF3 & Security &
    No plaintext credentials in source code. \texttt{xpack.security} disabled
    for development only — Keycloak integration is the planned production path. \\
  \rowcolor{rowalt}
  NF4 & Usability &
    Kibana dashboards are readable by a non-Elasticsearch expert. Severity is
    colour-coded: red (critical), orange (high), yellow (medium), green (info). \\
  NF5 & Reliability &
    Docker health-checks on Elasticsearch ensure Logstash and Filebeat only
    start after the database is accepting requests. \\
  \rowcolor{rowalt}
  NF6 & Portability &
    The entire stack (4 containers) launches with a single
    \texttt{docker compose up} command on any Docker-capable host. \\
  NF7 & Maintainability &
    Detection rules are defined in a single \field{RULES} dictionary in
    \file{alerts.py}. Adding a new rule requires adding one dict entry and
    one function — no schema changes. \\
  \bottomrule
  \caption{Non-Functional Requirements}
\end{longtable}

\subsection{Traceability Matrix}

\begin{center}
\begin{tabular}{p{4.5cm}p{4cm}p{5cm}}
  \toprule
  \textbf{Requirement} & \textbf{Component} & \textbf{File} \\
  \midrule
  F1 — Log Collection         & Filebeat            & \file{filebeat/filebeat.yml} \\
  F2/F3/F4 — Parse \& Enrich  & Logstash            & \file{logstash/pipeline/siem.conf} \\
  F5/F6 — Alert Generation    & Alert Engine        & \file{alerts.py} \\
  F7 — Blocklist              & Alert Engine        & \file{logs/blocklist.json} \\
  F8 — Dashboard              & Kibana 8.13         & Index pattern \field{siem-logs-*} \\
  F9 — Audit Trail            & Elasticsearch       & Index \field{siem-logs-*} \\
  F10 — Alert Index           & Logstash + ES       & Index \field{siem-alerts-*} \\
  NF5 — Reliability           & Docker Compose      & \file{docker-compose.yml} \\
  NF7 — Maintainability       & Alert Engine        & \file{alerts.py} \texttt{RULES} dict \\
  \bottomrule
\end{tabular}
\end{center}

\newpage

% ══════════════════════════════════════════════════════════════════════════════
\section{System Design}

\subsection{High-Level Architecture}

The system follows a \textbf{unidirectional pipeline architecture}. Data flows
through discrete, loosely-coupled stages. No component has knowledge of the
component two steps ahead of it, making each stage independently replaceable.

\begin{center}
\begin{tcolorbox}[colback=lightblue,colframe=accent,width=\linewidth,
  title={\textbf{End-to-End Data Flow}},fonttitle=\bfseries]
\small\centering
\begin{tabular}{c}
  \textbf{[1] Data Generation} \\
  \file{logg\_gen.py} writes SSH log lines to \file{logs/a.log} \\
  \file{alerts.py} writes JSON alerts to \file{logs/alerts.log} \\[0.5em]
  $\Big\downarrow$ \\[0.3em]
  \textbf{[2] Collection — Filebeat (port 5044)} \\
  Tails both log files; tags \field{log\_type: auth} or \field{log\_type: alert} \\[0.5em]
  $\Big\downarrow$ \\[0.3em]
  \textbf{[3] Processing — Logstash} \\
  Grok parse $\rightarrow$ Classify $\rightarrow$ Enrich $\rightarrow$ Timestamp $\rightarrow$ Route \\[0.5em]
  $\Big\downarrow$ \hspace{4cm} $\Big\downarrow$ \\[0.2em]
  \textbf{siem-logs-YYYY.MM.dd} \hspace{2cm} \textbf{siem-alerts-YYYY.MM.dd} \\
  \textit{(Elasticsearch)} \hspace{3.2cm} \textit{(Elasticsearch)} \\[0.5em]
  $\Big\uparrow$ polls every 10s \hspace{2.5cm} $\Big\uparrow$ re-ingested by Filebeat \\[0.2em]
  \textbf{[4] Detection — alerts.py} $\xrightarrow{\text{writes}}$ \file{logs/alerts.log} \\
  Also writes $\rightarrow$ \file{logs/blocklist.json} \\[0.5em]
  $\Big\downarrow$ \\[0.3em]
  \textbf{[5] Visualisation — Kibana (port 5601)} \\
  SSH Logs Overview dashboard \quad Alerts \& Threats dashboard \\
\end{tabular}
\end{tcolorbox}
\end{center}

\textbf{Why a pipeline architecture?} Each stage communicates only through
well-defined interfaces (a log file, a TCP port, an Elasticsearch index). This
means:
\begin{itemize}[leftmargin=2em,topsep=2pt]
  \item Logstash can be swapped for OpenSearch Ingest without touching
        Filebeat or the alert engine.
  \item The alert engine is a pure Python process — no framework coupling —
        making it trivial to unit test or replace with a different language.
  \item The \textbf{feedback loop} (alerts re-entering the pipeline as
        structured documents) is elegant: the same Logstash pipeline that
        handles raw logs also handles JSON alerts, keeping Kibana unified.
\end{itemize}

\subsection{Component Descriptions}

\begin{center}
\begin{tabular}{|p{3.2cm}|p{2.8cm}|p{7.5cm}|}
  \hline
  \rowcolor{tabhead}
  \textcolor{white}{\textbf{Layer}} &
  \textcolor{white}{\textbf{Component}} &
  \textcolor{white}{\textbf{Responsibility}} \\
  \hline
  Data Generation   & \file{logg\_gen.py}   & Simulates realistic SSH traffic with fixed IPs, random attacker pools, burst events, forced alert scenarios \\
  \hline
  \rowcolor{rowalt}
  Data Generation   & \file{log\_generator.py} & Legacy simple generator (1 line/3s, 70\% fail) — superseded by \file{logg\_gen.py} \\
  \hline
  Collection        & Filebeat 8.13         & Tails log files; ships lines to Logstash; maintains read-position registry across restarts \\
  \hline
  \rowcolor{rowalt}
  Processing        & Logstash 8.13         & Grok parsing, conditional classification (14 branches), timestamp normalisation, index routing \\
  \hline
  Storage           & Elasticsearch 8.13    & Inverted-index document store; daily time-based indices; provides aggregation API for alert engine \\
  \hline
  \rowcolor{rowalt}
  Detection         & \file{alerts.py}      & Polls Elasticsearch every 10s; evaluates 5 rules; writes JSON alerts; maintains IP blocklist \\
  \hline
  Visualisation     & Kibana 8.13           & Two dashboards with 14 panels; Discover for ad-hoc queries; index pattern management \\
  \hline
  Orchestration     & Docker Compose        & Defines all 4 containers, their networks, volumes, health-checks, and startup ordering \\
  \hline
\end{tabular}
\end{center}

\subsection{Data Model}

\subsubsection*{siem-logs-* Elasticsearch Document Schema}

\begin{center}
\begin{tabular}{p{3.5cm}p{1.8cm}p{7.8cm}}
  \toprule
  \textbf{Field} & \textbf{Type} & \textbf{Description / Example} \\
  \midrule
  \field{@timestamp}      & date    & Reconstructed from log fields; normalised to UTC (source: IST) \\
  \field{event\_type}     & keyword & \texttt{failed\_password}, \texttt{accepted\_publickey}, \texttt{sudo\_session} \ldots \\
  \field{severity}        & keyword & \texttt{info}, \texttt{warning}, \texttt{high}, \texttt{critical} \\
  \field{src\_ip}         & keyword & Source IP of the connection attempt \\
  \field{username}        & keyword & Targeted account name \\
  \field{auth\_method}    & keyword & \texttt{password} or \texttt{publickey} \\
  \field{login\_result}   & keyword & \texttt{success} or \texttt{failure} \\
  \field{user\_validity}  & keyword & \texttt{valid} or \texttt{invalid} (for failed logins) \\
  \field{server}          & keyword & Hostname of the target server (e.g.\ \texttt{server}) \\
  \field{service}         & keyword & Source service: \texttt{sshd}, \texttt{sudo}, \texttt{CRON}, \texttt{systemd-logind} \\
  \field{pid}             & integer & Process ID of the logging service \\
  \field{src\_port}       & integer & Ephemeral source port of the connection \\
  \field{key\_type}       & keyword & RSA / ECDSA / ED25519 (publickey logins only) \\
  \field{key\_fingerprint}& keyword & SHA256 fingerprint (publickey logins only) \\
  \field{log\_type}       & keyword & Always \texttt{auth} for raw SSH logs \\
  \bottomrule
\end{tabular}
\end{center}

\subsubsection*{siem-alerts-* Elasticsearch Document Schema}

\begin{center}
\begin{tabular}{p{3.5cm}p{1.8cm}p{7.8cm}}
  \toprule
  \textbf{Field} & \textbf{Type} & \textbf{Description / Example} \\
  \midrule
  \field{@timestamp}        & date    & UTC time the alert was fired by the engine \\
  \field{rule}              & keyword & \texttt{BRUTE\_FORCE}, \texttt{BURST\_ATTACK}, \texttt{ROOT\_LOGIN} \ldots \\
  \field{severity}          & keyword & \texttt{MEDIUM}, \texttt{HIGH}, \texttt{CRITICAL} \\
  \field{src\_ip}           & keyword & IP that triggered the rule \\
  \field{message}           & text    & Human-readable explanation (e.g.\ ``15 failed logins in 60s'') \\
  \field{failed\_attempts}  & integer & Count of events (brute force / burst rules) \\
  \field{window\_sec}       & integer & Detection window in seconds \\
  \field{targeted\_users}   & keyword & Usernames targeted in the attack window \\
  \field{log\_type}         & keyword & Always \texttt{alert} \\
  \bottomrule
\end{tabular}
\end{center}

\subsection{Design Trade-offs}

\begin{center}
\begin{tabular}{p{3cm}p{4.5cm}p{4cm}p{2cm}}
  \toprule
  \textbf{Decision} & \textbf{Choice Made} & \textbf{Alternative} & \textbf{Why} \\
  \midrule
  Log shipper      & Filebeat (lightweight, binary) & Fluentd (more plugins) & Lower memory; purpose-built for Elastic \\
  Storage          & Elasticsearch (aggregations) & PostgreSQL & No SQL aggregation pipeline; ES is purpose-built for log analytics \\
  Alert engine     & Python daemon & Kibana Watcher & Watcher uses complex YAML/JSON rule DSL; Python is testable and readable \\
  Alert feedback   & Re-ingest via Filebeat & Separate index writer & Keeps the pipeline uniform; no extra code path \\
  Index strategy   & Daily indices \texttt{siem-logs-YYYY.MM.dd} & Single index & Easy data retention; old indices deleted without rebuild \\
  Auth             & None (dev only) & Keycloak RBAC & Scope: demonstrates detection pipeline, not auth management \\
  \bottomrule
\end{tabular}
\end{center}

\newpage

% ══════════════════════════════════════════════════════════════════════════════
\section{Implementation Details}

\subsection{Tech Stack Justification}

\begin{itemize}[leftmargin=2em]
  \item \textbf{Python 3} — Alert engine and log generator. The Elasticsearch
        Python client provides a clean API for aggregation queries; the
        standard library covers JSON, datetime, and pathlib without extra
        dependencies.
  \item \textbf{Elastic Stack 8.13 (ELK)} — Industry-standard SIEM backbone.
        Grok, aggregations, Kibana Lens, and the Beats ecosystem are
        purpose-built for log analytics at scale. The 8.x series added native
        security features and improved Kibana dashboarding.
  \item \textbf{Filebeat} — Ultra-lightweight log shipper (written in Go).
        Maintains a read-position registry so no lines are skipped or
        duplicated across restarts.
  \item \textbf{Docker Compose} — Single-command reproducible environment.
        Health-check \texttt{depends\_on} enforces Elasticsearch readiness
        before dependent services start, eliminating startup race conditions.
\end{itemize}

\subsection{Module: logg\_gen.py — Realistic Traffic Simulator}

This is the most sophisticated module in the project. It produces SSH log
traffic that is indistinguishable from a real server under attack, ensuring
the entire detection pipeline always has meaningful data to process.

\subsubsection*{Fixed IP Scheduler}
Seven known actors are defined with type, schedule interval range, and
success rate:

\begin{center}
\begin{tabular}{p{3.2cm}p{1.8cm}p{2.8cm}p{2cm}p{3cm}}
  \toprule
  \textbf{IP} & \textbf{Type} & \textbf{Interval (s)} & \textbf{Success \%} & \textbf{Behaviour} \\
  \midrule
  \texttt{10.0.1.25}       & internal & 60--180   & 92\% & Corp workstation \\
  \texttt{192.168.1.105}   & internal & 90--240   & 88\% & Corp workstation \\
  \texttt{172.16.0.55}     & vpn      & 120--300  & 85\% & VPN remote user \\
  \texttt{34.201.12.45}    & legit    & 180--480  & 80\% & Cloud deploy host \\
  \texttt{185.220.101.12}  & tor      & 30--90    & 4\%  & Persistent Tor attacker \\
  \texttt{45.155.205.233}  & botnet   & 20--60    & 3\%  & Botnet node \\
  \texttt{103.214.132.55}  & scanner  & 15--45    & 2\%  & Port/credential scanner \\
  \bottomrule
\end{tabular}
\end{center}

A \texttt{FixedIPScheduler} class stores \texttt{next\_fire} timestamps and
re-schedules each IP after firing, randomised within its interval band.

\subsubsection*{Random IP Pools}
At startup, 40 attacker IPs and 10 legit IPs are generated from a function
that avoids RFC1918, loopback, and reserved ranges. Fresh IPs every session
prevent simple pattern matching from defeating the detector.

\subsubsection*{Burst State Machine}
A \texttt{BurstState} object tracks active burst campaigns. On each tick,
a 1.5\% random chance triggers a burst of 20--60 failed logins from one
attacker IP in rapid succession. The burst flag pauses all other event
generation until the campaign is complete.

\subsubsection*{Forced Alert Scenarios}
Every 10 seconds, two scenarios from a rotating list fire deterministically,
guaranteeing that the alert engine always has threshold-crossing events to
detect regardless of random traffic patterns:

\begin{center}
\begin{tabular}{p{4cm}p{5cm}p{4.5cm}}
  \toprule
  \textbf{Scenario} & \textbf{What It Does} & \textbf{Rule It Triggers} \\
  \midrule
  \texttt{scenario\_brute\_force}  & 15 failed logins from one IP in 1.5s   & \texttt{BRUTE\_FORCE} \\
  \texttt{scenario\_burst}         & 35 failed logins from \texttt{99.99.99.99} in 2s & \texttt{BURST\_ATTACK} \\
  \texttt{scenario\_root\_login}   & \texttt{Accepted password for root}      & \texttt{ROOT\_LOGIN} \\
  \texttt{scenario\_preauth\_storm}& 6 preauth disconnects from one IP in 0.6s & \texttt{PREAUTH\_STORM} \\
  \texttt{scenario\_max\_auth}     & \texttt{max authentication attempts exceeded} & \texttt{MAX\_AUTH\_EXCEEDED} \\
  \bottomrule
\end{tabular}
\end{center}

\subsubsection*{14 Log Line Builders}
Each function produces a syntactically correct syslog-format string with
the custom \texttt{YYYY Mon DD HH:MM:SS} timestamp prefix required by the
Logstash date filter:

\begin{lstlisting}[language=Python, caption=Sample log line builders from logg\_gen.py]
def line_accepted_password(user, ip, port):
    return (f"{ts()} server sshd[{next_pid()}]: "
            f"Accepted password for {user} from {ip} port {port} ssh2")

def line_failed(user, ip, port, invalid=False):
    tag = "invalid user " if invalid else ""
    return (f"{ts()} server sshd[{next_pid()}]: "
            f"Failed password for {tag}{user} from {ip} port {port} ssh2")

def line_accepted_pubkey(user, ip, port):
    kt = random.choice(["RSA", "ECDSA", "ED25519"])
    return (f"{ts()} server sshd[{next_pid()}]: Accepted publickey for "
            f"{user} from {ip} port {port} ssh2: {kt} SHA256:{sha256()}")
\end{lstlisting}

\subsection{Module: logstash/pipeline/siem.conf — Parsing Pipeline}

The Logstash pipeline is the structural core of the system. It transforms
free-text log lines into richly-typed Elasticsearch documents.

\subsubsection*{Stage 1 — Alert bypass}
If \field{log\_type == alert}, the document is a pre-parsed JSON object
from the alert engine. Logstash skips all Grok processing and routes it
directly to \field{siem-alerts-*}.

\subsubsection*{Stage 2 — Base Grok parse}
Two patterns cover the syslog envelope (with and without PID bracket):

\begin{lstlisting}[language=bash, caption=Base Grok patterns in siem.conf]
# With PID: "2026 Apr 10 14:22:05 server sshd[10234]: ..."
"^%{YEAR:log_year} %{MONTH:log_month} %{MONTHDAY:log_day}
 %{TIME:log_time} %{HOSTNAME:server}
 (?<service>[^\[\s]+)\[%{NUMBER:pid}\]: %{GREEDYDATA:log_message}$"

# Without PID: "2026 Apr 10 14:22:05 server sudo: ..."
"^%{YEAR:log_year} %{MONTH:log_month} %{MONTHDAY:log_day}
 %{TIME:log_time} %{HOSTNAME:server}
 (?<service>[^\[\s:]+): %{GREEDYDATA:log_message}$"
\end{lstlisting}

\subsubsection*{Stage 3 — 14-branch event classification}

\begin{center}
\begin{tabular}{p{4.5cm}p{2cm}p{6.5cm}}
  \toprule
  \textbf{event\_type} & \textbf{Severity} & \textbf{Match Pattern} \\
  \midrule
  \texttt{accepted\_password}    & info/critical* & \texttt{Accepted password for ...} \\
  \texttt{accepted\_publickey}   & info/critical* & \texttt{Accepted publickey for ...} \\
  \texttt{failed\_password} (invalid) & warning  & \texttt{Failed password for invalid user ...} \\
  \texttt{failed\_password} (valid)   & high     & \texttt{Failed password for ...} \\
  \texttt{max\_auth\_exceeded}   & high           & \texttt{maximum authentication attempts exceeded} \\
  \texttt{pam\_failure}          & warning        & \texttt{PAM N more authentication failure} \\
  \texttt{preauth\_disconnect}   & warning        & \texttt{Disconnected from invalid user ... [preauth]} \\
  \texttt{session\_disconnect}   & info           & \texttt{Disconnected from user ...} \\
  \texttt{connection\_closed}    & info           & \texttt{Connection closed by ...} \\
  \texttt{sudo\_session}         & info           & \texttt{pam\_unix(sudo:session)} \\
  \texttt{session\_opened}       & info           & \texttt{New session for user} \\
  \texttt{session\_closed}       & info           & \texttt{Session removed for user} \\
  \texttt{cron\_job}             & info           & \texttt{service == CRON} \\
  \texttt{other}                 & info           & Parsed envelope, unclassified message \\
  \bottomrule
  \multicolumn{3}{l}{\small *critical when \texttt{username == root}}
\end{tabular}
\end{center}

\subsubsection*{Stage 4 — Timestamp reconstruction}
SSH logs carry no year field. The pipeline concatenates extracted date
parts and parses them with the \texttt{date} filter in IST timezone,
overwriting \field{@timestamp} with the correct UTC value:

\begin{lstlisting}[language=bash, caption=Logstash timestamp handling]
mutate {
  add_field => { "log_timestamp" => "%{log_year} %{log_month}
                  %{log_day} %{log_time}" }
}
date {
  match    => ["log_timestamp", "yyyy MMM dd HH:mm:ss"]
  target   => "@timestamp"
  timezone => "Asia/Kolkata"
}
\end{lstlisting}

\subsubsection*{Stage 5 — Routing}
The \field{log\_type} field set by Filebeat routes output:
\begin{itemize}[leftmargin=2em,topsep=2pt]
  \item \field{log\_type: auth} $\rightarrow$ \field{siem-logs-YYYY.MM.dd}
  \item \field{log\_type: alert} $\rightarrow$ \field{siem-alerts-YYYY.MM.dd}
\end{itemize}

\subsection{Module: alerts.py — Detection Engine}

\begin{rulebox}[Detection Rules Configuration]
\begin{lstlisting}[language=Python]
RULES = {
    "BRUTE_FORCE":       {"failed_in_window": 10, "window_sec": 60,  "severity": "HIGH"},
    "BURST_ATTACK":      {"failed_in_window": 30, "window_sec": 60,  "severity": "CRITICAL"},
    "ROOT_LOGIN":        {"severity": "CRITICAL"},
    "MAX_AUTH_EXCEEDED": {"severity": "HIGH"},
    "PREAUTH_STORM":     {"count": 5, "window_sec": 60, "severity": "MEDIUM"},
}
POLL_INTERVAL = 10   # seconds between full rule evaluation
LOOKBACK      = 120  # seconds of history for single-event rules
COOLDOWN      = 60   # seconds before same rule fires again for same IP
\end{lstlisting}
\end{rulebox}

\subsubsection*{Rule Evaluation Flow}
Every 10 seconds the main loop calls all five rule functions in sequence.
Each function issues an Elasticsearch aggregation query, checks results
against its threshold, verifies the per-IP cooldown has expired, then
calls \texttt{write\_alert()} if the condition is met.

\begin{center}
\begin{tabular}{p{3.8cm}p{5.5cm}p{2.2cm}p{2cm}}
  \toprule
  \textbf{Rule} & \textbf{Elasticsearch Query} & \textbf{Threshold} & \textbf{Action} \\
  \midrule
  BRUTE\_FORCE       & Agg: failed logins by IP in 60s    & $\geq 10$ / IP & Alert + Block \\
  BURST\_ATTACK      & Agg: failed logins by IP in 60s    & $\geq 30$ / IP & Alert + Block \\
  ROOT\_LOGIN        & Filter: accepted + username=root   & Any            & Alert \\
  MAX\_AUTH\_EXCEEDED & Filter: max\_auth\_exceeded events & Any            & Alert \\
  PREAUTH\_STORM     & Agg: preauth disconnects by IP     & $\geq 5$ / IP  & Alert \\
  \bottomrule
\end{tabular}
\end{center}

\subsubsection*{Alert Document Structure}
Each alert is written as a single JSON line to \file{logs/alerts.log}:

\begin{lstlisting}[language=Python, caption=Sample alert document]
{
  "timestamp":       "2026-04-10T09:15:32.441Z",
  "rule":            "BRUTE_FORCE",
  "severity":        "HIGH",
  "src_ip":          "6.6.6.6",
  "failed_attempts": 15,
  "window_sec":      60,
  "targeted_users":  ["admin", "root"],
  "message":         "15 failed logins in 60s"
}
\end{lstlisting}

\subsubsection*{IP Blocklist}
\file{blocklist.json} accumulates blocked IPs across sessions:

\begin{lstlisting}[language=Python, caption=Sample blocklist.json]
{
  "6.6.6.6": {
    "blocked_at": "2026-04-10T09:15:32Z",
    "reason":     "Brute force: 15 attempts in 60s",
    "severity":   "HIGH"
  }
}
\end{lstlisting}

\subsection{Repository Structure}

\begin{lstlisting}[language=bash, caption=Project file tree]
siem-project/
├── docker-compose.yml          # 4-container infrastructure definition
├── logg_gen.py                 # Realistic SSH traffic simulator (primary)
├── log_generator.py            # Legacy simple generator (superseded)
├── alerts.py                   # Python detection engine + blocklist
├── logs/
│   ├── a.log                   # SSH log sink  (Filebeat input 1)
│   ├── alerts.log              # Alert JSON sink (Filebeat input 2)
│   └── blocklist.json          # Auto-blocked IPs (persistent)
├── filebeat/
│   └── filebeat.yml            # Two inputs + Logstash output
└── logstash/
    ├── config/logstash.yml     # Logstash settings
    └── pipeline/siem.conf      # Full Grok pipeline (14 event types)
\end{lstlisting}

\newpage

% ══════════════════════════════════════════════════════════════════════════════
\section{Kibana Dashboard Design}

Both dashboards use the wildcard index patterns \field{siem-logs-*} and
\field{siem-alerts-*} so they work across all daily indices without
reconfiguration. The default time range is set to \textbf{Last 24 hours}
and stored with each dashboard.

\subsection{Dashboard 1 — SSH Logs Overview}

\begin{center}
\begin{tabular}{p{4cm}p{2.5cm}p{2cm}p{4.5cm}}
  \toprule
  \textbf{Panel} & \textbf{Chart Type} & \textbf{Index} & \textbf{Key Config} \\
  \midrule
  Total Login Events       & Metric      & siem-logs-* & Count \\
  Total Failed Logins      & Metric      & siem-logs-* & Filter: \field{login\_result: failure} \\
  Total Successful Logins  & Metric      & siem-logs-* & Filter: \field{login\_result: success} \\
  Unique Attacker IPs      & Metric      & siem-logs-* & Unique count of \field{src\_ip.keyword} \\
  Login Activity Over Time & Line chart  & siem-logs-* & X: \field{@timestamp}; split by \field{login\_result} \\
  Top 10 Attacking IPs     & Bar (vert.) & siem-logs-* & Top values \field{src\_ip.keyword} size 10 \\
  Top 10 Targeted Usernames& Bar (vert.) & siem-logs-* & Top values \field{username.keyword} size 10 \\
  Event Type Distribution  & Donut       & siem-logs-* & Slice by \field{event\_type.keyword} size 10 \\
  Severity Over Time       & Bar (stack) & siem-logs-* & Split by \field{severity.keyword} \\
  Auth Method Breakdown    & Pie         & siem-logs-* & Slice by \field{auth\_method.keyword} \\
  Audit Trail              & Table       & siem-logs-* & Cols: timestamp, ip, user, event, severity \\
  \bottomrule
\end{tabular}
\end{center}

\textbf{Layout:}
\begin{lstlisting}[language=bash]
Row 1: [Total Events] [Failed] [Successful] [Unique Attacker IPs]
Row 2: [          Login Activity Over Time (full width)          ]
Row 3: [Top 10 IPs (half width)]  [Top 10 Usernames (half width)]
Row 4: [Event Type Donut (third)] [Auth Method Pie] [Severity Time]
Row 5: [          Audit Trail Table (full width)                 ]
\end{lstlisting}

\subsection{Dashboard 2 — Alerts \& Threats}

\begin{center}
\begin{tabular}{p{4cm}p{2.5cm}p{2.5cm}p{4cm}}
  \toprule
  \textbf{Panel} & \textbf{Chart Type} & \textbf{Index} & \textbf{Key Config} \\
  \midrule
  Total Alerts             & Metric       & siem-alerts-* & Count \\
  Critical Alerts          & Metric (red) & siem-alerts-* & Filter: \field{severity: CRITICAL} \\
  High Alerts              & Metric (org.)& siem-alerts-* & Filter: \field{severity: HIGH} \\
  Unique Threat IPs        & Metric       & siem-alerts-* & Unique count \field{src\_ip.keyword} \\
  Alert Timeline           & Line chart   & siem-alerts-* & Split by \field{severity.keyword} \\
  Alerts by Rule           & Bar (horiz.) & siem-alerts-* & Y: \field{rule.keyword}, X: count \\
  Alert Severity Dist.     & Donut        & siem-alerts-* & Slice by \field{severity.keyword} \\
  Top 10 Threat IPs        & Bar (horiz.) & siem-alerts-* & Y: \field{src\_ip.keyword} size 10 \\
  Live Alert Feed          & Table        & siem-alerts-* & Cols: time, rule, severity, ip, message \\
  \bottomrule
\end{tabular}
\end{center}

\textbf{Layout:}
\begin{lstlisting}[language=bash]
Row 1: [Total Alerts] [Critical] [High] [Unique Threat IPs]
Row 2: [          Alert Timeline (full width)              ]
Row 3: [Alerts by Rule (half)]  [Severity Donut (half)    ]
Row 4: [Top 10 Threat IPs (half)]  [                      ]
Row 5: [          Live Alert Feed Table (full width)       ]
\end{lstlisting}

\newpage

% ══════════════════════════════════════════════════════════════════════════════
\section{Security Considerations}

\begin{infobox}[Development vs.\ Production Security Posture]
The current deployment is a development/demonstration environment.
The considerations below describe the gap between the current state and
a production-hardened deployment.
\end{infobox}

\subsection{Current Security Measures}

\begin{itemize}[leftmargin=2em]
  \item \textbf{No credentials in source code} — Elasticsearch URL is
        passed via environment variable \texttt{ES\_URL}; no passwords are
        hard-coded.
  \item \textbf{Network isolation} — All containers communicate on a private
        Docker bridge network (\texttt{siem-net}). Only ports 9200, 5601, and
        5044 are exposed to localhost; no public exposure.
  \item \textbf{Read-only mounts} — Filebeat configuration and log directories
        are mounted as \texttt{:ro} (read-only) in the containers where write
        access is not required.
  \item \textbf{Automatic threat blocking} — HIGH and CRITICAL alerts
        auto-populate the blocklist. In a production deployment this file would
        feed an \texttt{iptables} / firewall rule set.
\end{itemize}

\subsection{Planned Production Hardening}

\begin{center}
\begin{tabular}{p{4cm}p{9cm}}
  \toprule
  \textbf{Area} & \textbf{Planned Measure} \\
  \midrule
  Authentication     & Keycloak OIDC for Admin / Analyst / Read-only roles; JWT bearer tokens on all API paths \\
  Elasticsearch      & Enable \texttt{xpack.security}; TLS on all inter-container communication; RBAC per index \\
  Secrets management & Docker Secrets or HashiCorp Vault for Elasticsearch credentials \\
  OWASP ZAP          & Automated DAST scan on the Kibana surface and any API endpoints exposed externally \\
  Log integrity      & WORM-style write-once log storage; cryptographic hash chaining on \file{a.log} \\
  Blocklist enforcement & \texttt{cron} job that reads \file{blocklist.json} and inserts \texttt{iptables DROP} rules \\
  \bottomrule
\end{tabular}
\end{center}

\newpage

% ══════════════════════════════════════════════════════════════════════════════
\section{Testing \& Validation}

\subsection{Testing Strategy}

\begin{center}
\begin{tabular}{p{3.5cm}p{3cm}p{6.5cm}}
  \toprule
  \textbf{Type} & \textbf{Tool} & \textbf{Scope} \\
  \midrule
  Pipeline integration & Docker Compose + Kibana Discover & Full end-to-end smoke test from log write to ES document \\
  Alert rule validation & Forced scenarios in \file{logg\_gen.py} & All 5 detection rules triggered deterministically \\
  Grok pattern testing  & Kibana Grok Debugger & Each of the 14 log format variants \\
  Manual API testing    & \texttt{curl} / Postman & Elasticsearch \texttt{\_search} and \texttt{\_count} endpoints \\
  Security scan         & OWASP ZAP & Kibana web surface \\
  UI automation         & Selenium WebDriver & Dashboard load, metric card values \\
  Code quality          & SonarQube & Python modules (\file{alerts.py}, \file{logg\_gen.py}) \\
  \bottomrule
\end{tabular}
\end{center}

\subsection{Sample Test Cases}

\begin{longtable}{p{0.6cm}p{4cm}p{5cm}p{4cm}}
  \toprule
  \textbf{TC} & \textbf{Test} & \textbf{Steps} & \textbf{Expected} \\
  \midrule\endhead
  T1 & Grok parses failed login &
    Append one \texttt{Failed password for invalid user} line to \file{a.log} &
    Document in ES: \field{event\_type=failed\_password}, \field{severity=warning} \\
  \midrule
  T2 & Grok parses publickey login &
    Append \texttt{Accepted publickey for alice ... ED25519 SHA256:...} &
    Fields \field{key\_type=ED25519} and \field{key\_fingerprint} present \\
  \midrule
  T3 & BRUTE\_FORCE fires &
    Run \texttt{scenario\_brute\_force()} (15 attempts, 1.5s) &
    Alert in \field{siem-alerts-*} within 20s; \field{rule=BRUTE\_FORCE} \\
  \midrule
  T4 & ROOT\_LOGIN fires &
    Append \texttt{Accepted password for root from 1.2.3.4} &
    CRITICAL alert within 15s; IP added to \file{blocklist.json} \\
  \midrule
  T5 & Cooldown suppresses duplicate &
    Fire same scenario twice within 30s &
    Exactly one alert entry per IP per rule; second suppressed \\
  \midrule
  T6 & BURST\_ATTACK fires &
    Run \texttt{scenario\_burst()} (35 attempts, 2s) &
    Alert with \field{severity=CRITICAL}, \field{rule=BURST\_ATTACK} \\
  \midrule
  T7 & Blocklist auto-populated &
    Trigger any HIGH rule &
    IP entry appears in \file{blocklist.json} with timestamp and reason \\
  \midrule
  T8 & Kibana dashboard loads &
    Open \texttt{localhost:5601}, navigate to SSH Logs Overview &
    All 11 panels render with data; no ``No results'' panels \\
  \midrule
  T9 & Wildcard index covers all days &
    Run generator across midnight &
    Both \field{siem-logs-YYYY.MM.DD} and next-day index appear in dashboard \\
  \midrule
  T10 & Alert pipeline round-trip &
    Trigger PREAUTH\_STORM; check Alerts dashboard &
    Alert visible in \field{siem-alerts-*}; Live Alert Feed table shows entry \\
  \bottomrule
  \caption{Sample Test Cases}
\end{longtable}

\subsection{Bugs Found \& Fixed}

\begin{center}
\begin{tabular}{p{0.7cm}p{5cm}p{7.5cm}}
  \toprule
  \textbf{ID} & \textbf{Bug} & \textbf{Resolution} \\
  \midrule
  B1 & Grok \texttt{\_grokparsefailure} on all lines &
    Service field character class \texttt{[\^{}[ ]+} was too greedy and matched
    the bracket. Refined to \texttt{[\^{}[\textbackslash{}s]+} to stop at bracket or whitespace. \\
  B2 & Filebeat silently fails to start &
    Bind-mounted \file{filebeat.yml} triggered strict ownership check.
    Added \texttt{filebeat -e -strict.perms=false} to Compose command. \\
  B3 & Kibana shows no data after new day &
    Dashboard was bound to daily index \texttt{siem-logs-2026.04.05}.
    Replaced with wildcard \texttt{siem-logs-*} pattern. \\
  B4 & Thousands of duplicate alerts per minute &
    Every 10s poll re-fired rules for IPs still in the lookback window.
    Added per-IP 60s cooldown dict \texttt{\_fired} in \file{alerts.py}. \\
  B5 & \field{@timestamp} set to ingest time, not log time &
    Logstash was not overwriting \field{@timestamp}. Added \texttt{date} filter
    with \texttt{target => ``@timestamp''} and IST timezone. \\
  \bottomrule
\end{tabular}
\end{center}

\newpage

% ══════════════════════════════════════════════════════════════════════════════
\section{Demo \& Screencast}

\subsection{Screencast Link}

\begin{infobox}[Screencast]
\textbf{Video Link:} \textit{[Insert your MP4 / YouTube / Google Drive link here]}\\
\textbf{Duration:} $\approx$10 minutes \quad \textbf{Format:} MP4 with voiceover
\end{infobox}

\subsection{Timestamped Demo Script}

\begin{enumerate}[leftmargin=2em]
  \item \textbf{[0:00--1:00] Architecture overview}\\
    Show the architecture flow diagram. Walk through the repository structure
    in the terminal. Explain the unidirectional pipeline in one minute.
  \item \textbf{[1:00--2:30] Starting the stack}\\
    Run \texttt{docker compose up -d}. Show the health-check waiting loop for
    Elasticsearch. Confirm all four containers are healthy with
    \texttt{docker compose ps}. Open Kibana at \texttt{localhost:5601}.
  \item \textbf{[2:30--4:30] Live log generation}\\
    Run \texttt{python3 logg\_gen.py} in one terminal. Show the scrolling
    output with legitimate logins, failed attempts, and forced scenarios.
    Switch to Kibana Discover and show documents appearing in
    \field{siem-logs-*} in real time. Point out parsed fields.
  \item \textbf{[4:30--6:30] SSH Logs Overview dashboard}\\
    Open Dashboard 1. Walk through each panel: metric row, attack timeline
    (show the spike when a burst fires), top attacking IPs, event type donut,
    severity heatmap, and the audit trail table. Filter to critical only using
    the filter bar.
  \item \textbf{[6:30--8:30] Alerts \& Threats dashboard}\\
    In a second terminal run \texttt{python3 alerts.py}. Wait 10 seconds and
    switch to Dashboard 2. Show the alert metric cards increment. Click the
    Live Alert Feed table — point out BRUTE\_FORCE and ROOT\_LOGIN entries.
    Show the severity donut and ``Alerts by Rule'' bar chart filling up.
  \item \textbf{[8:30--9:30] Blocklist}\\
    Open \file{logs/blocklist.json} in the terminal with \texttt{cat}.
    Show auto-blocked IPs with timestamps and reasons. Explain how this
    would feed \texttt{iptables} in production.
  \item \textbf{[9:30--10:00] Wrap-up}\\
    Summarise what was demonstrated. Mention Keycloak RBAC, GeoIP enrichment,
    and ML anomaly detection as the next planned steps.
\end{enumerate}

\newpage

% ══════════════════════════════════════════════════════════════════════════════
\section{Lessons Learned \& Future Work}

\subsection{What Worked Well}

\begin{itemize}[leftmargin=2em]
  \item \textbf{ELK pipeline reliability} — Once the Grok patterns were
        correct, the pipeline required zero manual intervention. Logs flowed
        end-to-end without restarts, service interruptions, or data loss.
  \item \textbf{Feedback-loop architecture} — Writing alerts as JSON lines to
        a log file and re-ingesting them via Filebeat was an elegant design.
        The alert data became queryable in Kibana alongside raw logs without
        any extra tooling or a second Elasticsearch client in the alert engine.
  \item \textbf{Forced scenarios in the log generator} — Embedding
        deterministic alert-triggering events directly in \file{logg\_gen.py}
        made rule validation trivial. No separate test harness was needed.
  \item \textbf{Docker Compose health-checks} — The
        \texttt{condition: service\_healthy} dependency prevented hours of
        race-condition debugging. Every run starts cleanly in the correct
        order regardless of machine speed.
  \item \textbf{Daily index strategy} — Using \texttt{siem-logs-YYYY.MM.dd}
        indices made data retention simple (delete old indices without
        rebuilding) and gave the dashboard a natural time-based structure.
\end{itemize}

\subsection{What Failed / Challenges Faced}

\begin{itemize}[leftmargin=2em]
  \item \textbf{Grok pattern iteration} — Getting all 14 SSH log variants
        parsed correctly required over twenty iterations in the Kibana Grok
        Debugger. The most problematic: services with hyphens
        (\texttt{systemd-logind}) and PAM lines with semicolon-separated
        key=value pairs.
  \item \textbf{Year in timestamp} — Standard SSH syslog format omits the
        year. We added a year prefix in \file{logg\_gen.py} and built a
        matching Logstash date pattern. Testing across midnight revealed an
        edge case where the year rolled forward in the log but the date filter
        used the previous year.
  \item \textbf{Filebeat strict permissions} — The bind-mounted
        \file{filebeat.yml} failed a silent ownership security check, causing
        Filebeat to start and immediately exit. The error was only visible via
        \texttt{docker logs siem-filebeat} and took significant time to
        diagnose.
  \item \textbf{Alert cooldown design} — Initial testing without a cooldown
        produced 500+ duplicate alerts per minute as the 10s poll cycle
        repeatedly saw the same IPs in the 120s lookback window. The
        in-memory \texttt{\_fired} dict with COOLDOWN=60 solved this but
        required careful testing to ensure legitimate new bursts from the same
        IP still triggered alerts after the cooldown expired.
  \item \textbf{Kibana index pattern scope} — Early dashboards were built
        against the daily index pattern, which required manual recreation every
        day. Switching to wildcard patterns late in development meant
        rebuilding several visualisations.
\end{itemize}

\subsection{Top 3 Future Features}

\begin{center}
\begin{tabular}{p{0.4cm}p{3.8cm}p{7.5cm}p{1.8cm}}
  \toprule
  \textbf{\#} & \textbf{Feature} & \textbf{Description} & \textbf{Effort} \\
  \midrule
  1 & Keycloak RBAC &
    Integrate Keycloak for Admin / Analyst / Read-only roles.
    Protect Kibana with OIDC. Admin-only access to blocklist management
    and alert acknowledgement.
    & 2 weeks \\
  \midrule
  2 & GeoIP Map Visualisation &
    Add the Logstash GeoIP filter to enrich \field{src\_ip} with country,
    city, and lat/lon coordinates. Display attacker origins as a heat map
    on Kibana Maps layer. Instantly shows geographic attack origin patterns.
    & 1 week \\
  \midrule
  3 & ML Anomaly Detection &
    Use Elasticsearch Machine Learning jobs to detect: unusual login times
    for a known user, impossible travel (same user, two IPs, short delta),
    and off-hours CRON executions. Feeds into a new ``Anomaly'' alert rule.
    & 3--4 weeks \\
  \bottomrule
\end{tabular}
\end{center}

\subsection{Hindsight — Tips for Future Teams}

\begin{enumerate}[leftmargin=2em]
  \item \textbf{Design your log format before writing any pipeline code.}
        Every field you want in Elasticsearch must be extractable by Grok.
        Changing the log format mid-project cascades into new patterns,
        re-indexing, and broken Kibana mappings simultaneously.
  \item \textbf{Use the Kibana Grok Debugger obsessively.}
        It gives immediate feedback on pattern matches without restarting a
        Logstash container. Every pattern in \file{siem.conf} was developed
        and validated here first.
  \item \textbf{Create wildcard index patterns on day one.}
        \texttt{siem-logs-*} instead of \texttt{siem-logs-2026.04.10}.
        A daily pattern means rebuilding your dashboards every single day
        you run the project.
  \item \textbf{Read container logs before assuming the issue is in your code.}
        Three of our five bugs (B2, B3, B4) were only visible in
        \texttt{docker logs}. \texttt{docker logs -f siem-filebeat} and
        \texttt{docker logs -f siem-logstash} should be your first debugging
        tools.
\end{enumerate}

\subsection{AI Integration Opportunities}

\begin{itemize}[leftmargin=2em]
  \item \textbf{Dynamic alert thresholds} — An ML model trained on historical
        login patterns could set BRUTE\_FORCE thresholds per-server and
        per-time-of-day, eliminating false positives during legitimate
        high-traffic periods (e.g., Monday morning login bursts).
  \item \textbf{Natural language log search} — An LLM interface over the
        Elasticsearch API would let an admin query:
        \textit{``Show me all root logins from non-internal IPs last week''}
        without knowing Kibana KQL syntax.
  \item \textbf{Automated incident summaries} — When an alert fires, an LLM
        could generate a one-paragraph incident report: who, what IP, how many
        attempts, usernames targeted, similar past incidents, recommended
        response.
  \item \textbf{Grok pattern generation} — During development, an AI assistant
        could generate and iterate Grok patterns from sample log lines,
        replacing the most time-consuming part of the project.
\end{itemize}

\newpage

% ══════════════════════════════════════════════════════════════════════════════
\section{Team Contribution Statement}

\begin{infobox}[Note to Markers]
Replace the placeholder names and contribution details below with your actual
team members before submission. The commit history on GitHub provides a
verifiable audit trail of individual contributions.
\end{infobox}

\begin{longtable}{p{2.8cm}p{4cm}p{6.5cm}}
  \toprule
  \textbf{Member} & \textbf{Primary Area} & \textbf{Contributions} \\
  \midrule\endhead
  \textit{[Name 1]} & Infrastructure \& Pipeline &
    \file{docker-compose.yml} setup, Filebeat configuration, Logstash Grok
    pipeline design (all 14 patterns), Elasticsearch index management,
    Docker health-check configuration. \\
  \midrule
  \textit{[Name 2]} & Log Generator &
    \file{logg\_gen.py}: fixed IP scheduler, random pools, burst state machine,
    forced alert scenarios, all 14 log line builder functions.
    \file{log\_generator.py} (initial prototype). \\
  \midrule
  \textit{[Name 3]} & Alert Engine \& Blocklist &
    \file{alerts.py}: all 5 detection rules, cooldown logic, blocklist writer,
    alert JSON schema design. Tuned \texttt{POLL\_INTERVAL},
    \texttt{LOOKBACK}, and \texttt{COOLDOWN} values. \\
  \midrule
  \textit{[Name 4]} & Dashboards \& Testing &
    Kibana SSH Logs Overview dashboard (11 panels), Alerts \& Threats dashboard
    (9 panels), index pattern management, Selenium test scripts, OWASP ZAP
    security scan, SonarQube code quality review. \\
  \bottomrule
  \caption{Team Contributions}
\end{longtable}

\subsection*{Coordination Approach}

\begin{itemize}[leftmargin=2em]
  \item Weekly standups via \textit{[platform]} — 30-minute sessions to
        review what was completed, what is blocked, and what is next.
  \item \textbf{GitHub Issues} used for bug tracking with assignees, labels
        (\texttt{bug}, \texttt{enhancement}, \texttt{pipeline}), and
        milestone assignments.
  \item \textbf{Feature branches} merged to \texttt{main} via pull requests
        requiring at least one team-member approval before merge.
  \item Commit messages follow the format
        \texttt{[component] short description} for easy filtering in
        \texttt{git log}.
\end{itemize}

\newpage

% ══════════════════════════════════════════════════════════════════════════════
\section{Appendix}

\subsection{A — Full Detection Rules Reference}

\begin{center}
\begin{tabular}{p{4cm}p{1.8cm}p{2.8cm}p{2cm}p{3cm}}
  \toprule
  \textbf{Rule} & \textbf{Window} & \textbf{Threshold} & \textbf{Severity} & \textbf{Auto Action} \\
  \midrule
  BRUTE\_FORCE        & 60s  & $\geq 10$ failed/IP   & HIGH     & Alert + Block IP \\
  BURST\_ATTACK       & 60s  & $\geq 30$ failed/IP   & CRITICAL & Alert + Block IP \\
  ROOT\_LOGIN         & 120s & Any root acceptance   & CRITICAL & Alert only \\
  MAX\_AUTH\_EXCEEDED & 120s & Any occurrence        & HIGH     & Alert only \\
  PREAUTH\_STORM      & 60s  & $\geq 5$ preauth/IP   & MEDIUM   & Alert only \\
  \bottomrule
\end{tabular}
\end{center}

\subsection{B — Filebeat Input Configuration}

\begin{lstlisting}[language=bash, caption=filebeat/filebeat.yml]
filebeat.inputs:
  - type: log
    enabled: true
    paths: [/logs/a.log]
    scan_frequency: 5s
    close_inactive: 10s
    fields:
      log_type: auth
    fields_under_root: true

  - type: log
    enabled: true
    paths: [/logs/alerts.log]
    scan_frequency: 5s
    close_inactive: 10s
    json.keys_under_root: true
    json.add_error_key: true
    fields:
      log_type: alert
    fields_under_root: true

output.logstash:
  hosts: ["logstash:5044"]
\end{lstlisting}

\subsection{C — Docker Compose Service Summary}

\begin{center}
\begin{tabular}{p{3.5cm}p{2cm}p{3cm}p{4.5cm}}
  \toprule
  \textbf{Container} & \textbf{Port} & \textbf{Image} & \textbf{Health Check} \\
  \midrule
  siem-elasticsearch & 9200  & elasticsearch:8.13.0 & \texttt{curl \_cluster/health} \\
  siem-kibana        & 5601  & kibana:8.13.0        & Depends on ES healthy \\
  siem-logstash      & 5044  & logstash:8.13.0      & Depends on ES healthy \\
  siem-filebeat      & —     & filebeat:8.13.0      & Depends on logstash \\
  \bottomrule
\end{tabular}
\end{center}

\subsection{D — Kibana Dashboard Panels Summary}

\begin{center}
\small
\begin{tabular}{p{3cm}p{4cm}p{2.5cm}p{3.5cm}}
  \toprule
  \textbf{Dashboard} & \textbf{Panel Name} & \textbf{Type} & \textbf{Index} \\
  \midrule
  Logs Overview & Total / Failed / Successful / Unique IPs & Metric ×4 & siem-logs-* \\
  Logs Overview & Login Activity Over Time & Line & siem-logs-* \\
  Logs Overview & Top 10 Attacking IPs & Bar (vert.) & siem-logs-* \\
  Logs Overview & Top 10 Targeted Usernames & Bar (vert.) & siem-logs-* \\
  Logs Overview & Event Type Distribution & Donut & siem-logs-* \\
  Logs Overview & Severity Over Time & Bar (stack) & siem-logs-* \\
  Logs Overview & Auth Method Breakdown & Pie & siem-logs-* \\
  Logs Overview & Audit Trail & Table & siem-logs-* \\
  Alerts \& Threats & Total / Critical / High / Unique IPs & Metric ×4 & siem-alerts-* \\
  Alerts \& Threats & Alert Timeline & Line & siem-alerts-* \\
  Alerts \& Threats & Alerts by Rule & Bar (horiz.) & siem-alerts-* \\
  Alerts \& Threats & Alert Severity Distribution & Donut & siem-alerts-* \\
  Alerts \& Threats & Top 10 Threat IPs & Bar (horiz.) & siem-alerts-* \\
  Alerts \& Threats & Live Alert Feed & Table & siem-alerts-* \\
  \bottomrule
\end{tabular}
\end{center}

\subsection{E — Installation Quick Reference}

\begin{lstlisting}[language=bash, caption=Full stack startup commands]
# Prerequisites: Docker Desktop running, Python 3 installed
# Recommended: 8 GB RAM allocated to Docker

# Step 1 — Start the ELK stack (wait ~60s for Elasticsearch)
cd siem-project
docker compose up -d
docker compose ps   # verify all containers healthy

# Step 2 — Generate realistic SSH log traffic
python3 logg_gen.py

# Step 3 — Start the alert detection engine
python3 alerts.py

# Step 4 — Open Kibana
open http://localhost:5601
# Navigate to: Stack Management > Index Patterns
# Create: siem-logs-*   (time field: @timestamp)
# Create: siem-alerts-* (time field: @timestamp)
# Then: Dashboards > Create dashboard > follow Section 5 design

# Useful debugging commands
docker logs -f siem-filebeat    # watch Filebeat shipping lines
docker logs -f siem-logstash    # watch Logstash parsing output
cat logs/blocklist.json         # view auto-blocked IPs
\end{lstlisting}

\vfill
\begin{center}
  \textcolor{accent}{\rule{0.6\linewidth}{0.4pt}}\\[0.6em]
  {\small\textcolor{gray}{SIEM Project — Software Engineering Submission — \today}}
\end{center}

\end{document}
